\section{Conclusion}

This project investigated the accuracy of machine learning regression models in predicting the
number of GitHub stargazers for open-source repositories. We collected a dataset of 1{,}000
repositories via the GitHub REST API, extracted 14 activity-based features, and trained five
regression models under a consistent evaluation protocol using the $R^2$ metric.

Our results show that \textbf{[BEST MODEL NAME]} achieves the highest predictive accuracy
($R^2 = \text{XX.XX}$), outperforming both linear baselines and the other ensemble methods
considered. The result confirms that non-linear models are better suited to capturing the
complex relationships between repository activity signals and community popularity.

The prediction system was deployed to a production environment using a fully automated pipeline.
OpenStack provisioned the virtual machines, Ansible configured the infrastructure, and a Git
Hook eliminated the need for any manual deployment steps. The production Flask application,
backed by Celery and Redis, demonstrated horizontal scalability: throughput increased by
approximately \textbf{XX\%} when the number of Celery workers was scaled from one to four.

Future work could explore richer feature sets---such as commit frequency, contributor diversity,
and documentation quality---as well as time-series models that capture the temporal dynamics of
star accumulation. Extending the dataset beyond 1{,}000 repositories and incorporating
cross-platform signals (e.g.\ social media mentions) could further improve prediction accuracy.
